\scanpage{4}
\mspage{3}

\textcircled{\scriptsize c}\quad $\forall\,T\in\mathrm{Mod}_{f}(\unsure{\mathcal L|L})$, et $A\in C(T)/\unsure{R|\mathcal R}_T$
une classe suivant $R_T$, une partie
\unsure{de} $\dot A\subset T$, appelée \uline{frontière} de $A$.

\unsure{Ces|Les} données \textcircled{\scriptsize a} \textcircled{\scriptsize b}, \textcircled{\scriptsize c} étant \unsure{assujetties}
à satisfaire certains axiomes.

Je prévois que les axiomes impliquent
\deleted{Soit \illegible{1} $\unsure{t|U}\in C(T)$, \unsure{on pose}}
ceci

1°) que la donnée \textcircled{\scriptsize c} résulte déjà de \textcircled{\scriptsize b}

2°) que la donnée \textcircled{\scriptsize b} résulte de
la donnée des $R_T$, \deleted{\unsure{pour}} dans le cas
particulier où $T\in S_0$. Donc
les seules \unsure{donnée|données} essentielles \unsure{serait|seraient}
\textcircled{\scriptsize a} \unsure{celles|celle} des 0-sphères i.e. la notion
de \deleted{\unsure{pos.}} lieux en pos. relative modérée
(\unsure{relation|condition} \inserted{réflexive et} symétrique) et \textcircled{\scriptsize b} La donnée,
pour \deleted{\unsure{tout}} $\unsure{\varepsilon|\epsilon|E}\in S_0$, de la relation
d'équivalence $R_{\varepsilon}$ \unsure{dans} $C(\varepsilon)$.

\deleted{\unsure{On va|Donc on}} \uline{Axiomes provisoires :}

\deleted{Soit \unsure{$X$}} \unsure{Appellons|Appelons} les $U$ de la forme
$C(\deleted{\unsure{L}}T)$ \uline{parties modérées cofinies}. On

% paper watermark
\rotatebox{90}{\textsc{\unsure{S}trong Extra St}}
